\documentclass{beamer}

\mode<presentation> {%
	\usetheme[]{Boadilla}
	\usecolortheme{orchid}
	\setbeamercovered{transparent}
	\setbeamersize{text margin left=0.7cm}
  \setbeamersize{text margin right=0.7cm}
  \setbeamertemplate{enumerate item}[square]
  \setbeamertemplate{enumerate subitem}[circle]
  \setbeamertemplate{enumerate subsubitem}[>]
  \setbeamertemplate{itemize item}[square]
  \setbeamertemplate{itemize subitem}[circle]
  \setbeamertemplate{itemize subsubitem}[>]
}

% \usepackage{bm}
% \usepackage[utf8]{inputenc}
% \usepackage[spanish]{babel}
% \usepackage{times}
% \usepackage[T1]{fontenc}
% \usepackage{graphicx}
% \usepackage{multirow}
% \usepackage{pifont}
% \usepackage{rotating}
% \usepackage{colortbl}
% \usepackage[framed]{mcode}
\usepackage[utf8]{inputenc}
\usepackage{amsmath}
\usepackage{graphicx}
\usepackage{booktabs}

% \definecolor{olivegreen}{rgb}{0.0,0.5,0.0}
% \definecolor{general}{rgb}{0.9 0.9 1}%
% \newrgbcolor{colorsistema}{0 1 1}
% \newrgbcolor{colorviolet}{0.9 0.8 1}

\title[Procesamiento de Señales]{Transformada de Fourier Discreta (DFT)}
\author[J. Chacón]{%[Author, Another] % (optional, use only with lots of authors)
\textit{Ingeniería Electrónica de Comunicaciones }\\
\vspace{0.6cm}
\small Jesús Chacón Sombría}

\institute[DACYA] % (optional, but mostly needed)
{
  %\inst{1}%

\footnotesize  Departamento de Arquitectura de Computadores y Automática.\\
  Universidad Complutense de Madrid
  }

\date[Transformada Fourier Discreta] % (optional)
{\footnotesize Curso 2025-2026}%{\today / Master}
\subject{Talks}

\pgfdeclareimage[height=0.75cm]{logo}{ucm.png}

\usepackage[absolute,overlay]{textpos}
\setlength{\TPHorizModule}{1mm}
\setlength{\TPVertModule}{1mm}
\newcommand{\MyLogo}{%
\begin{textblock}{14}(116.0,0.6)
  \pgfuseimage{logo}
\end{textblock}
}

\makeatletter
\numberwithin{section}{part}
\AtBeginPart{
\beamer@tocsectionnumber=0\relax
\frame{\partpage}
}
\makeatother

\begin{document}

\begin{frame}
  \titlepage
\end{frame}

\newcommand*\oldmacro{}%
\let\oldmacro\insertshorttitle%
\renewcommand*\insertshorttitle{%
 \MyLogo
  \oldmacro\hfill%
}

\AtBeginSection[]{}

\AtBeginSubsection[]{}      

\section{Introducción}

\begin{frame}
  \frametitle{¿Qué es la Transformada Discreta de Fourier (DFT)?}
  \begin{itemize}
      \item La Transformada Discreta de Fourier (DFT) es una herramienta matemática fundamental en el procesamiento de señales.
      \item Convierte una secuencia de valores discretos en una secuencia de coeficientes en el dominio de la frecuencia.
      \item Es ampliamente utilizada en aplicaciones como análisis espectral, compresión de audio, y procesamiento de imágenes.
      \item Existen algoritmos eficientes para calcular la DFT.
  \end{itemize}    
\end{frame}    

\begin{frame}
  \frametitle{Transformada de Fourier de Tiempo Discreto (DTFT)}
  \begin{alertblock}{Transformada de Fourier de Tiempo Discreto (DTFT)}
  \[
    \begin{aligned}
      X(\omega) &=\substack{\infty\\\sum\\n=-\infty}x[n]\text{e}^{-j\omega n}\\
      x[n] &=\frac{1}{2\pi}\int_0^{2\pi}X(\omega)\text{e}^{j\omega n}d\omega
    \end{aligned}
  \]
  \end{alertblock}
  
  \begin{itemize}
    \item Antes de definir y utilizar la DFT, recordamos conceptos importantes.
    \item La DTFT $X(\omega)$ es una función:
    \begin{itemize} 
      \item continua en $\omega$, 
      \item periódica (con periodo $2\pi$).
    \end{itemize}
  \end{itemize}

\end{frame}

\begin{frame} 
  \frametitle{Relación entre espectro de señales continuas y discretas}
    \begin{center}
    \alert{¿Podemos conocer el espectro $X(\omega)$ de una señal $x(t)$\\  
    a partir de sus muestras $\color{blue}x_d[n]=x(nT_s)$?}
    \end{center}
    \begin{itemize}
      \item Sabemos que $\color{blue}X_d(\Omega=\omega T_s)=\frac{1}{T_s}X_a(\omega)$, con $\color{blue}\Omega=\omega T_s\in(-\pi,\pi)$.
      \item Pero conocemos un número finito de muestras: $\color{blue}\hat{x}[n]=x_d[n]w[n]\color{black}$ \vspace{0.2cm}
      \begin{itemize}
        \item $w[n]=\left\{ \begin{matrix}1~si~0\leq n< N\\ 0~c.c.\end{matrix} \right. \leftrightarrow W(\Omega)=\frac{1-e^{j\Omega N}}{1-e^{j\Omega}}$
        \item $\hat{X}(\Omega)=\int_{0}^{2\pi}X_d(\xi)W(\Omega-\xi)d\xi$ \\ \vspace{0.1cm}
        \alert{Truncar la señal en el tiempo distorsiona el espectro}
      \end{itemize} \vspace{0.2cm}  
      \item Estimación del espectro: $\color{blue}X_a(\omega)=T_s X_d(\omega)\approx T_s \hat{X_d}(\omega)\color{black}$ \vspace{0.2cm} 
      \item Estimación de la densidad espectral: $\color{blue}|\Phi_{x_a x_a}(\omega)|\approx T_s^2|\Phi_{\hat{x}\hat{x}}(\omega)|\color{black}$
    \end{itemize}
\end{frame}    

\begin{frame}
  \frametitle{Definición Matemática de la DFT}
  Al igual que en el dominio del tiempo, es útil muestrear la DTFT.
  \begin{itemize}
    \item Únicamente muestras de un periodo (habitualmente $\color{blue}\omega\in[0,2\pi]$).
    \item Si elegimos $N$ muestras, el periodo resultante es $\color{blue}\Delta\omega=2\pi/N$.
    \item Entonces, $\color{blue}X[k]=X(\frac{2\pi k}{N})$.
  \end{itemize}

  \begin{alertblock}{Transformada de Fourier Discreta (DFT)}
    La DFT de una secuencia \( x[n] \) de longitud \( N \) se define como:
    \[
    X[k] = \sum_{n=0}^{N-1} x[n] \cdot e^{-j \frac{2\pi}{N} kn} \quad \text{para } k = 0, 1, \dots, N-1  
    \]
    donde:
    \begin{itemize}
        \item \( X[k] \) es el coeficiente en el dominio de la frecuencia.
        \item \( x[n] \) es la secuencia en el dominio del tiempo.
        \item \( N \) es el número de puntos en la secuencia.
    \end{itemize}
  \end{alertblock}
\end{frame}    

\begin{frame}
  \alert {¿Podemos reconstruir $x[n]$ a partir de $X[k$]?}
  \begin{itemize}
    \item Partimos de la DTFT muestreada: \vspace{-0.2cm} \[ X[k]=X(\frac{2\pi k}{N})=\substack{\infty\\\sum\\n=-\infty}x[n]\text{e}^{-j\frac{2\pi k}{N} n} \]
    \item Separamos el sumatorio en grupos de $N$ muestras.
    \[
      \begin{aligned}
        X[k]&= ... +
        \substack{-1\\\sum\\n=-N}x[n]\text{e}^{-j\frac{2\pi k}{N} n} +
        \substack{N-1\\\sum\\n=0}x[n]\text{e}^{-j\frac{2\pi k}{N} n} +
        \substack{2N-1\\\sum\\n=N}x[n]\text{e}^{-j\frac{2\pi k}{N} n} + ...
        \\
          &= \substack{N-1\\\sum\\n=0}\left[\substack{\infty\\\sum\\m=-\infty}x[n-m N]\right]\text{e}^{-j\frac{2\pi k}{N} (n - mN)}
      \end{aligned}
    \]
    \item Definimos $\color{blue}x_p[n]=\left[\substack{\infty\\\sum\\m=-\infty}x[n-m N]\right]$.
    \item $x_p[n]$ es la extensión periódica de $x[n]$.
  \end{itemize}
\end{frame}

\begin{frame}
  \alert {¿Podemos reconstruir $x[n]$ a partir de $X[k$]?}
  \[X[k]=\substack{N-1\\\sum\\n=0}x_p[n]\text{e}^{-j\frac{2\pi k}{N} n }\]
  \begin{itemize}
    \item Es fácil comprobar que $x_p[n]$ es periódica, con periodo $N$.
    \item Llamamos a $x_p[n]$ \textbf{la extensión periódica} de $x[n]$.
    \item Para reconstruir $x[n]$ se debe cumplir $x_p[n]=x[n]$ para $0\leq n<L$
  \end{itemize}

  \begin{alertblock}{Condiciones para poder reconstruir $x[n]$}
    \begin{itemize}
      \item (i) $x[n]$ tiene longitud finita $L$,
      \item (ii) el número de muestras de la DFT es $N\geq L$.
    \end{itemize}
    \begin{center}
      $X[k]=\substack{N-1\\\sum\\n=0}x[n]\text{e}^{-j\frac{2\pi k}{N} n}$\\
      $x[k]=\frac{1}{N}\substack{N-1\\\sum\\n=0}X[k]\text{e}^{j\frac{2\pi k}{N} n}$
    \end{center}
  \end{alertblock}
\end{frame}

% \section{Propiedades de la DFT}

% \begin{frame}
%     \frametitle{Algunas propiedades de la DFT}
%     \begin{itemize}
%         \item \textbf{Linealidad}: \( \text{DFT}\{a x[n] + b y[n]\} = a X[k] + b Y[k] \)
%         \item \textbf{Desplazamiento en el tiempo}: \( \text{DFT}\{x[n - m]\} = X[k] e^{-j \frac{2\pi}{N} km} \)
%         \item \textbf{Simetría}: Si \( x[n] \) es real, entonces \( X[k] = X^*[N - k] \)
%         \item \textbf{Convolución}: \( \text{DFT}\{x[n] * y[n]\} = X[k] \cdot Y[k] \)
%     \end{itemize}    
% \end{frame}    

% \section{Aplicaciones}

% \begin{frame}
%     \frametitle{Aplicaciones de la DFT}
%     \begin{itemize}
%         \item \textbf{Análisis espectral}: Identificación de frecuencias en señales.
%         \item \textbf{Compresión de audio}: MP3, AAC, etc.
%         \item \textbf{Procesamiento de imágenes}: Filtrado, compresión (JPEG), etc.
%         \item \textbf{Comunicaciones digitales}: Modulación OFDM.
%     \end{itemize}    
% \end{frame}    

% \section{Ejemplo Práctico}

% \begin{frame}
%     \frametitle{Ejemplo Práctico}
%     \begin{itemize}
%         \item Consideremos una señal sinusoidal simple:
%         \[
%         x[n] = \cos\left(\frac{2\pi}{N} k_0 n\right)  
%         \]
%         \item Su DFT tendrá picos en \( k = k_0 \) y \( k = N - k_0 \).
%         \item \textbf{Ejemplo gráfico}:
%     \end{itemize}    
%     \begin{center}
%         % \includegraphics[width=0.6\textwidth]{dft_example.png}
%     \end{center}    
% \end{frame}    





\section{Análisis Espectral con la DFT}

\begin{frame}
    \frametitle{Análisis Espectral con la DFT}
  \begin{itemize}
    \item ¿Qué es el análisis espectral?
    \begin{itemize}
        \item Estudia cómo se distribuye la energía de una señal en frecuencia.
        \item La densidad espectral de potencia es una medida importante.
        \item La DFT se utiliza para estimar esta densidad espectral.
    \end{itemize}    
    \item La energía se conserva entre tiempo y frecuencia.
  \end{itemize}    
  \begin{alertblock}{\centering Teorema de Plancherel - Identidad de Parseval}
    \vspace{0.2cm} \centering
    $\int_{-\infty}^{\infty}x(t)x^*(t)dt = \frac{1}{2\pi}\int_{-\infty}^{\infty}X(\omega)X^*(\omega)d\omega$ \\ \vspace{0.3cm}
    $\sum_{n=0}^{N-1}x[n] x^{*}[n] = \frac{1}{N}\sum_{k=0}^{N-1} X[k] X^*[k]$ \\ \vspace{0.3cm}
  \end{alertblock}  
\end{frame}    

\begin{frame}
    \frametitle{Análisis Espectral con la DFT}
  \begin{itemize}
    \item {\small Para señales definidas en energía: $\color{blue}E_x = \int_{-\infty}^{\infty}|x(t)|^2dt<\infty\color{blue}$}
    \begin{alertblock}{\centering Densidad espectral de energía}
      \vspace{0.2cm}
      \centering $S_{xx}(\omega)=X(\omega)X^*(\omega)=|X(\omega)|^2$ \\ \vspace{0.2cm}
    \end{alertblock}  
    \item {\small Para señales definidas en potencia: $\color{blue}P_x = \substack{lim\\T\rightarrow\infty}\frac{1}{T}\int_{-T/2}^{T/2}|x(t)|^2dt<\infty\color{blue}$}
    \begin{alertblock}{\centering Densidad espectral de potencia}
      Teorema de Wiener-Khinchin \\
      \centering $\Phi_{xx}(\omega)=\int_{-\infty}^{\infty}\phi_{xx}(t)\text{e}^{-j\omega t}dt$ \\ \vspace{0.2cm}
      $\phi_{xx}(\tau)=\substack{lim\\T\rightarrow\infty}\int_{-T/2}^{T/2}x(t)x(t+\tau)d\tau$ (autocorrelación) \\ \vspace{0.2cm}
    \end{alertblock}  
    \end{itemize}
  
\end{frame}    


\section{Representación e Interpretación de la DFT}


% \begin{frame}\frametitle{Conservación de energía entre tiempo y frecuencia}
%   \begin{alertblock}{\centering Teorema de Plancherel - Identidad de Parseval}
%     \vspace{0.2cm} \centering
%     $\int_{-\infty}^{\infty}x(t)x^*(t)dt = \frac{1}{2\pi}\int_{-\infty}^{\infty}X(\omega)X^*(\omega)d\omega$ \\ \vspace{0.3cm}
%     $\sum_{n=0}^{N-1}x[n] x^{*}[n] = \frac{1}{N}\sum_{k=0}^{N-1} X[k] X^*[k]$ \\ \vspace{0.3cm}
%   \end{alertblock}  

%   \begin{itemize}
%     \item Transformada de Fourier Discreta (DFT): \\
%     \vspace{0.2cm}
%     $\color{blue}x[n]=\frac{1}{N}\sum_{k=0}^{N-1}X[k]e^{\:j\frac{2\pi}{N}k n}\color{black}
%       \leftrightarrow
%      \color{blue}X[k]=\sum_{n=0}^{N-1}x[n]e^{-j\frac{2\pi}{N}k n}\color{black}$ 
%     \\ \vspace{0.2cm} 
%     % {\footnotesize $\color{cyan}x^*[n]=\frac{1}{N}\sum_{k=0}^{N-1}X^*[k]e^{-j\frac{2\pi}{N}k n}\color{black}
%     %   \leftrightarrow
%     %  \color{cyan}X^*[-k]=\sum_{n=0}^{N-1}x^*[n]e^{-j\frac{2\pi}{N}k n}\color{black}$ 
%     % } 
%     \vspace{0.2cm}
%     \item Potencia finita: $\color{blue}P_x = \frac{1}{N}\sum_{n=0}^{N-1}x[n]x^*[n]\color{black}$ \\ \vspace{0.2cm}
%     % {\footnotesize
%     %   $\color{cyan}\sum_{n=0}^{N-1}x[n] \left[\frac{1}{N}\sum_{k=0}^{N-1}X^*[k]e^{-j\frac{2\pi}{N}k n}\right]
%     %       = \frac{1}{N} \sum_{k=0}^{N-1} \left[\sum_{n=0}^{N-1}x[n]e^{-j\frac{2\pi}{N}k n}\right] X^*[k]\color{black}$
%     % }      
%   \end{itemize}  

% \end{frame}  

















% \section{Conclusión}

% \begin{frame}
%     \frametitle{Conclusión}
%     \begin{itemize}
%         \item La DFT es una herramienta poderosa para el análisis de señales en el dominio de la frecuencia.
%         \item Sus aplicaciones son vastas y fundamentales en muchas áreas de la ingeniería y la ciencia.
%         \item Comprender la DFT es esencial para el procesamiento moderno de señales.
%     \end{itemize}    
% \end{frame}    

% \begin{frame}
%     \frametitle{Preguntas}
%     \begin{center}
%         \Huge ¿Preguntas?
%     \end{center}    
% \end{frame}    



% \section{Transformada Discreta de Fourier (DFT)}

% \begin{frame}
%     \frametitle{Transformada Discreta de Fourier (DFT)}
%     \begin{itemize}
%         \item La DFT es una herramienta fundamental en el procesamiento digital de señales.
%         \item Convierte una señal en el dominio del tiempo en una representación en el dominio de la frecuencia.
%         \item Se define como:
%         \[
%         X[k] = \sum_{n=0}^{N-1} x[n] \cdot e^{-j \frac{2\pi}{N} kn} \quad \text{para } k = 0, 1, \dots, N-1
%         \]
%         \item Donde \(X[k]\) es el coeficiente en el dominio de la frecuencia, \(x[n]\) es la señal en el dominio del tiempo, y \(N\) es el número de puntos en la secuencia.
%     \end{itemize}
% \end{frame}

% \section{Transformada Rápida de Fourier (FFT)}

% \begin{frame}
%     \frametitle{Transformada Rápida de Fourier (FFT)}
%     \begin{itemize}
%         \item La FFT es un algoritmo eficiente para calcular la DFT.
%         \item Reduce la complejidad computacional de \(O(N^2)\) a \(O(N \log N)\).
%         \item Utiliza el enfoque de "divide y vencerás" para dividir la secuencia en subsecuencias más pequeñas y combinar los resultados de manera recursiva.
%     \end{itemize}
% \end{frame}


% \section{Conclusión}

% \begin{frame}
%     \frametitle{Conclusión}
%     \begin{itemize}
%         \item La DFT y la FFT son herramientas fundamentales en el procesamiento digital de señales.
%         \item Permiten analizar las componentes frecuenciales de una señal discreta.
%         \item Las tareas propuestas en la práctica permiten aplicar los conceptos teóricos en un caso práctico de análisis de una señal de audio.
%     \end{itemize}
% \end{frame}

\end{document}